\chapter{Cervical Laminectomy}
\index{Cervical Laminectomy}

\section*{Definition and Overview}
Cervical laminectomy is a surgical procedure in which the lamina -- the bony arch at the back of a vertebra -- is removed to relieve pressure on the spinal cord or the nerve roots in the neck. It is most commonly performed for cervical spondylotic myelopathy, a condition in which age-related narrowing of the spinal canal (spinal stenosis) compresses the spinal cord and causes weakness, numbness, and loss of coordination.

The operation creates more space within the spinal canal, allowing the spinal cord and nerves to lie with less pressure. It may be performed alone or combined with fusion of the affected vertebrae to maintain stability of the neck. Cervical laminectomy is major spinal surgery, and it is considered when conservative treatments have not controlled symptoms and the spinal cord is under significant compression.

\section*{Epidemiology}
Degenerative disease of the cervical spine is very common, and cervical spondylotic myelopathy is the most common cause of spinal cord dysfunction in people over the age of 55. As populations age, the number of people with spinal canal narrowing and cord compression is growing.

Most people with neck degeneration do not require surgery, and only a minority of those with myelopathy proceed to laminectomy. The decision to operate is individual and depends on the severity of cord compression, the rate of progression, and the person's symptoms and general fitness. Surgical rates vary widely between countries and between individual surgeons.

\section*{Aetiology and Pathophysiology}
Cervical laminectomy is performed to relieve compression of the spinal cord or nerve roots within the cervical spine. The conditions that cause this compression include:

\begin{itemize}
    \item \textbf{Cervical spondylosis:} age-related wear and tear of the discs and joints, with bone spurs (osteophytes) encroaching on the canal
    \item \textbf{Disc herniation:} protrusion of a cervical disc that presses on the spinal cord or nerve roots
    \item \textbf{Ligamentum flavum hypertrophy:} thickening of the ligament at the back of the canal narrowing the space
    \item \textbf{Congenital spinal stenosis:} a naturally narrow spinal canal that becomes symptomatic with age
    \item \textbf{Tumours:} growths within or around the spinal canal that compress the cord
    \item \textbf{Trauma:} fractures or dislocations of the cervical spine
    \item \textbf{Infection:} abscesses or inflammatory collections within the spinal canal
\end{itemize}

Removing the laminae opens the back of the canal and creates more room for the neural structures. In some cases, the surgeon also removes osteophytes or disc fragments through the same exposure.

\section*{Clinical Features}
Symptoms relate to compression of the spinal cord and nerve roots:

\begin{itemize}
    \item Neck pain and stiffness
    \item Numbness or tingling in the hands and fingers
    \item Weakness of the arms and hands, with clumsiness and dropping objects
    \item Loss of fine motor control, such as difficulty writing or doing up buttons
    \item Unsteadiness on the feet and a tendency to trip (spastic gait)
    \item A feeling of electric-shock-like sensations down the back or arms
    \item Changes in bladder or bowel function in advanced compression
    \item Pain radiating into the arms or shoulders if nerve roots are affected
\end{itemize}

\section*{Differential Diagnosis}
\begin{itemize}
    \item Cervical disc herniation without significant cord compression, which can often be managed conservatively
    \item Peripheral nerve disorders, such as carpal tunnel syndrome, which cause similar hand symptoms
    \item Multiple sclerosis and other demyelinating diseases of the spinal cord
    \item Motor neurone disease, which causes progressive weakness
    \item Vitamin B12 deficiency, which can cause cord-like symptoms
    \item Amyotrophic lateral sclerosis and other degenerative conditions of the cord
    \item Tumours or infections of the spine, which need specific investigation
\end{itemize}

\section*{When to Seek Medical Care}
See a doctor or spine specialist if you have persistent neck pain with numbness, weakness, or coordination problems in the hands, or increasing unsteadiness on your feet. Symptoms that are worsening or that affect daily function warrant prompt specialist assessment.

\textbf{Contact your local emergency services for:} sudden severe weakness in the arms or legs, loss of bladder or bowel control, or new difficulty walking, which may indicate an urgent spinal problem requiring immediate attention.

\section*{Diagnosis and Investigations}
\begin{itemize}
    \item \textbf{History and examination:} assessment of neck symptoms, hand function, gait, reflexes, and sensation
    \item \textbf{MRI of the cervical spine:} the key investigation; it shows the degree of cord compression and its cause
    \item \textbf{CT scan:} provides detailed bony detail, particularly useful when surgery is planned
    \item \textbf{X-rays:} may show alignment, instability, or advanced degeneration
    \item \textbf{Nerve conduction studies:} help exclude peripheral nerve conditions such as carpal tunnel syndrome
    \item \textbf{Pre-operative assessment:} general health checks, including heart, lungs, and blood tests, before surgery
\end{itemize}

\section*{Management}
\begin{itemize}
    \item \textbf{Conservative treatment:} physiotherapy, pain relief, and activity modification for milder symptoms
    \item \textbf{Cervical laminectomy:} removal of the laminae to decompress the spinal cord
    \item \textbf{Laminoplasty:} an alternative in which the canal is expanded by hinging the lamina open rather than removing it
    \item \textbf{Fusion:} combined with the laminectomy when instability is present, using bone graft and metal implants
    \item \textbf{Anterior surgery:} some compression is better treated from the front of the neck (anterior cervical decompression and fusion) rather than by laminectomy
    \item \textbf{Post-operative care:} wound care, pain management, early mobilisation with a cervical collar, and physiotherapy
    \item \textbf{Rehabilitation:} gradual return to activities with strengthening and gait training
\end{itemize}

\section*{Complications}
\begin{itemize}
    \item Infection of the wound, the disc space, or the bone, which can be serious
    \item Bleeding or a haematoma compressing the spinal cord
    \item Spinal fluid leak from a tear in the dura (the membrane around the cord)
    \item Nerve or spinal cord injury, causing new weakness or numbness
    \item Instability or deformity of the neck after removal of bone, sometimes requiring later fusion
    \item Failure of fusion if the procedure included a bone graft
    \item Persistent pain or incomplete relief of pre-operative symptoms
    \item Risks of general anaesthesia, including blood clots and chest infection
\end{itemize}

\section*{Prognosis and Outcomes}
Most people gain significant relief from the pain, numbness, and weakness that brought them to surgery, and improvement in walking and hand function commonly follows successful decompression. Outcomes are better when surgery is performed before severe or long-standing spinal cord damage has occurred. Recovery takes time: hospital stay is usually short, but full rehabilitation can continue for months. A minority of people have residual symptoms, and some will need further surgery for degeneration at other levels.

\section*{Prevention and Self-Care}
\begin{itemize}
    \item Maintain neck strength and good posture, and avoid heavy strain on the neck
    \item Keep a healthy weight and remain physically active to support the spine
    \item Do not smoke, as smoking impairs bone healing and disc health
    \item Report new numbness, weakness, or bladder symptoms promptly
    \item After surgery, follow activity restrictions and wear the cervical collar as advised
    \item Attend physiotherapy and follow-up appointments to maximise recovery
    \item Avoid heavy lifting and high-risk activities until cleared by the surgical team
\end{itemize}

\section*{Sources and Further Reading}
The following organisations provide reliable information on cervical spine surgery and spinal cord compression.

\begin{itemize}
    \item National Health Service. \textit{NHS conditions guide on cervical spine conditions and surgery.} https://www.nhs.uk/conditions/
    \item Mayo Clinic. \textit{Cervical laminectomy and spine surgery information.} https://www.mayoclinic.org/
    \item MSD Manual. \textit{Cervical spondylotic myelopathy -- professional overview.} https://www.msdmanuals.com/
    \item MedlinePlus. \textit{Consumer health information on spinal cord compression and laminectomy.} https://medlineplus.gov/
    \item National Institute for Health and Care Excellence. \textit{Clinical guidance on spinal surgery.} https://www.nice.org.uk/
    \item World Health Organization. \textit{Global guidance on rehabilitation after surgery and disability.} https://www.who.int/
\end{itemize}
